\chapter{Derivations for WP2}
\label{chapter:derivations-wp2}

\section{Analytical approximation of the solar-sail lightness number as a function of the shift in the Lagrange points.}

For an ideal solar sail oriented normal to the Sun-line, the required lightness number $\beta$ is defined by (\parencite{mcinnes1999book}):

$$ \beta = - \frac{1}{1-\mu} \frac{dU}{dx} $$

where $\mu$ is the mass ratio and the potential gradient along the synodic $x$-axis is \parencite{mcinnes1999book}:

$$ \frac{dU}{dx} = -\frac{(1-\mu)(x+\mu)}{r_1^3} - \frac{\mu(x-(1-\mu))}{r_2^3} + x $$

We perform a first-order Taylor series expansion of the potential gradient about the $L_2$ equilibrium point. By definition:

$$ \left. \frac{dU}{dx} \right|_{x_{L2}} = 0 $$

Evaluating the gradient at the shifted position $x = x_{L2} - \delta$ yields the first-order approximation:

$$ \frac{dU}{dx}\bigg|_{x_{L2} - \delta} \approx \underbrace{\frac{dU}{dx}\bigg|_{x_{L2}}}_{=0} - \delta \cdot \frac{d^2U}{dx^2}\bigg|_{x_{L2}} + \mathcal{O}(\delta^2) $$

The expression for the required lightness number simplifies to:

$$ \beta(\delta) \approx \frac{\delta}{1-\mu} \frac{d^2U}{dx^2} $$

The second derivative of the potential characterizes the local curvature of the potential field. Differentiating the gradient with respect to $x$:

$$ \frac{d^2U}{dx^2} = \frac{d}{dx} \left( -\frac{(1-\mu)(x+\mu)}{r_1^3} - \frac{\mu(x-(1-\mu))}{r_2^3} + x \right) $$

Given that $r_1 = |x + \mu|$ and $r_2 = |x - (1-\mu)|$ for collinear points ($y=z=0$), the expression for $\frac{d^2U}{dx^2}$ is:

$$ \frac{d^2U}{dx^2} = -\frac{(1-\mu)}{r_1^3} + \frac{3(1-\mu)(x+\mu)^2}{r_1^5} - \frac{\mu}{r_2^3} + \frac{3\mu(x-(1-\mu))^2}{r_2^5} + 1 $$

At the classical $L_2$ point, $(x+\mu)^2 = r_1^2$ and $(x-(1-\mu))^2 = r_2^2$. Consequently, the terms simplify to:

$$ \frac{d^2U}{dx^2} = \left( -\frac{1-\mu}{r_1^3} + \frac{3(1-\mu)}{r_1^3} \right) + \left( -\frac{\mu}{r_2^3} + \frac{3\mu}{r_2^3} \right) + 1 $$

$$ \frac{d^2U}{dx^2} = \frac{2(1-\mu)}{r_1^3} + \frac{2\mu}{r_2^3} + 1 $$

Substituting the simplified $\frac{d^2U}{dx^2}$ into the linearized equilibrium condition, we obtain the final analytical approximation for $\beta$ as a function of the displacement $\delta$:

$$ \boxed{ \beta(\delta) = \frac{\delta}{1-\mu} \left( \frac{2(1-\mu)}{r_1^3} + \frac{2\mu}{r_2^3} + 1 \right) } $$

This linear model establishes that the required sail lightness number is directly proportional to the magnitude of the shift from the classical equilibrium. The validity of this approximation is subject to two primary constraints:
\begin{itemize}
    \item The displacement $\delta$ must be small relative to the characteristic length scale of the local potential gradient, typically $\delta \ll r_2 \approx (\mu/3)^{1/3}$.
    \item The sail acceleration vector must remain strictly collinear with the Sun-planet line to satisfy the one-dimensional assumption.
\end{itemize}